\chapter{LA ENCUESTA}
\section{\emph{Propósito}}
\subsection{} Menéndez Pidal fue quien primero notó la característica más interesante del habla de los Montes de Pas\footnote{Véase `Pasiegos y vaqueiros'...}. En una excursión a estos valles verificó que allí las vocales finales [-u] e [-i] inflexionaban la vocal tónica y sugirió a los investigadores del ALPI que cuando hicieran sus encuestas en la provincia de Santander prestaran especial atención a este fenómeno, conocido desde hacía tiempo como característico del bable centro-meridional\footnote{Cf. M. Pidal, \emph{Notas acerca del bable de Lena}, 1899, y \emph{El dialecto leonés}, 1906.}. Uno de estos investigadores, el Sr. L. Rodríguez Castellano, publicó\footnote{En `Algunas precisiones sobre la metafonía de Santander y Asturias', 1959.} después los materiales que de este modo se habían recogido respecto a la metafonía. Fue este mismo erudito quien nos propuso el estudio más detenido del habla de una zona tan interesante.

No existía monografía de ningún habla de todo el territorio santanderino, y lo único que se había publicado en materia de dialectología montañesa eran unas cortas colecciones de voces y el vocabulario, en sus dos ediciones, de García Lomas (cf. la Bibliografía). En su libro \emph{El dialecto leonés}, M. Pidal hace muchas referencias al dialecto montañés, pero las demás obras que tratan de problemas dialectológicos más generales se limitan a emplear los datos publicados en las obras citadas. 

Dentro de la zona estudiada están comprendidos los puntos 407 (Vega de Pas) y 408 (Bustantegua) del ALPI. Dentro de la zona de metafonía (cf. mapa 4) están también comprendidos los puntos 411 (Resconorio) y 409 (Veguilla). Está muy cercano el punto 402 (Miera). 

Era de desear que se hiciera alguna monografía sobre el habla de una localidad del territorio montañés, zona tan arcaica donde se mezclan desde antiguo influjos castellanos y leoneses y que carecía de todo estudio sistemático. Nuestro propósito ha sido el de reparar en algo esta falta. 

\section{\emph{La comarca estudiada}}

\subsection{} Como ya se ha indicado (§ 11), hoy es preciso, para oír el verdadero lenguaje local, apartarse bastante de las carreteras. En todos los casos hemos empezado por preguntar en las cabezas de municipio por los puntos de habla más arcaizante y éstos han sido casi siempre caseríos sin comunicación fácil. Sin embargo, debemos insistir en que es inútil todo esfuerzo para caracterizar con alguna exactitud el habla de una localidad pasiega específica, ya que pocos pasiegos viven en un punto fijo y la gran mayoría vive sucesivamente en varios barrios y, muchas veces, aun en distintos municipios, dentro del territorio pasiego (cf. §§ 5-6).

Lo único que se puede notar, lo cual está indicado arriba (§ 4), es cierta diferenciación entre una subzona norte y otra sur (cf. mapa 3). Los fenómenos distintivos son mínimos; por ejemplo, en la subzona norte la -a final tiene un timbre más marcadamente palatal \textipa{[\H{5}]} o \textipa{[\`{@}]} que en el sur \textipa{[\'{5}]} (cf. § 23); hay mayor tendencia en el norte a ensordecer la vocal final que en el sur (cf. § 24); en la subzona sur las personas Nos y Vos del presente de los verbos en \textit{-er} e \textit{-ir} presentan las desinencias únicas [-émus -ís] ([salémus \textipa{ku\=rís}]), mientras en el norte se mantiene la distinción de conjugaciones (cf. § 128); finalmente, en el estudio de Cosas y Palabras (§§ 171-512) se notan numerosas variaciones léxicas entre las dos subzonas. Mientras en la subzona norte hay más arcaísmo, la subzona sur está más abierta a los influjos castellano-vulgares. 

Con las reservas arriba dichas, nuestro estudio se ha localizado en los puntos siguientes (cf. mapa 3): 


\begin{enumerate}[a)]
    \item \begin{tabular}{ p{3.5cm} l }
    municipio de Vega de Pas. &
    $\left\{
    \begin{tabular}{p{7cm}}
    
    \begin{enumerate}[1)]
       \item barrio de Pandillo, agrupación de casas más grande que lo corriente y muy arcaizante; no tiene carretera y está a 4 kms. de la cabeza del municipio.
       \item barrio de Yera, también a 4kms. de Vega de Pas; está a unos 500 m. de la carretera Vega de Pas-Espinosa. 
    \end{enumerate}
    \end{tabular}
    \right.$ \\
    \end{tabular}
    
    \item \begin{tabular}{ p{3.5cm} l }
    municipio de San Pedro del Romeral. &
    $\left\{
    \begin{tabular}{p{7cm}}
    
    \begin{enumerate}[1)]
       \item barrio de Bustaleguín, muy arcaizante, a 3 kms. de la carretera. 
       \item barrio de la Pradera, a medio kilómetro de San Pedro.  
    \end{enumerate}
    \end{tabular}
    \right.$ \\
    \end{tabular}

    \item 
    \begin{savenotes}
    \begin{tabular}{ p{3.5cm} l }
    municipio de Selaya. &
    $\left\{
    \begin{savenotes}
    \begin{tabular}{p{7cm}}
    
    \begin{enumerate}[1)]
       \item barrio de Cubilla (pertenece administrativamente al barrio de Bustantegua\footnote{Localidad notada como muy arcaizante por L. Rodriguez Castellano, `Algunas Precisiones...'}), pequeño caserío de unos 15-20 vecinos, muy arcaizante, a 5 kms. del pueblo de Selaya; desde hace unos años un ramal de la carretera llega hasta Valvanuz, a 2 kilómetros de Cubilla.
       \item barrio de Pisueña, a 7 kms. de Selaya; tiene carretera desde hace sólo diez años.  
    \end{enumerate}
    \end{tabular}
    \end{savenotes}
    \right.$
    \end{tabular}
    \end{savenotes}


    \item 
    \begin{savenotes}
    \begin{tabular}{ p{3.5cm} l }
    municipio de San Roque de Riomiera. &
    $\left\{
    \begin{savenotes}
    \begin{tabular}{p{7cm}}
    
    \begin{enumerate}[1)]
       \item barrio de Veolamadera (fon. \textipa{b\adp\ucs} < VADU) a 3 kms. del pueblo de San Roque y a una altura de 800 m. (400 m. sobre San Roque). 
       \item barrio de Calseca, a 3 kms. de San Roque\footnote{Calseca pertenece en realidad al municipio de Ruesga y está en la orilla derecha del Miera, pero lo hemos estudiado y lo incluimos aquí, puesto que está a 22 kms. de la cabeza de su propio municipio y toda su vida vierte hacia San Roque, a 3 kms. de distancia. No tiene carretera.}.    
    \end{enumerate}
    \end{tabular}
    \end{savenotes}
    \right.$ \\
    \end{tabular}
    \end{savenotes}
\end{enumerate}


\section{\emph{Condiciones y métodos de la encuesta.}}

\subsection{} El método principal de la encuesta ha sido el de la convivencia durante más de un año y medio con los habitantes de las localidades nombradas (§ 13). Durante ese período ha sido nuestra costumbre pasar cada semana dos o tres días, por lo menos, en los Montes de Pas, observando a los pasiegos en las faenas del campo y haciendo innumerables preguntas. Estas preguntas se han hecho a base de cuestionario (cf. abajo, § 15) y de conversación dirigida. Lo fundamental ha sido el cuestionario, pero el empleo de este instrumento no se ha limitado al interrogatorio directo, sino que ha servido de constante punto de referencia para la conversación dirigida. 

Para recoger los nombres de los objetos de uso diario ha bastado indicarlos, siempre que esto era posible, con el dedo, mientras para otros objetos hemos tenido que recurrir a la descripción. Dentro de lo posible, hemos procurado en los interrogatorios directos e indirectos, emplear el lenguaje de los mismos sujetos. 

A pesar del largo período de la investigación, éste no ha sido lo bastante prolongado para que por el método indirecto pudiéramos recoger en toda su riqueza las formas verbales, y para esta labor también hemos tenido que recorrer el cuestionario (§ 15). El método seguido ha sido el de mencionar un infinitivo, que sabíamos de antemano era dialectal, y después repetir una frase dialectal incompleta omitiendo el verbo, que iría en el último lugar. Indicábamos la persona y el tiempo por medio de pronombres y adverbios y el sujeto debía terminar en cada caso la frase. Variando los pronombres y adverbios, hemos podido completar los distintos paradigmas de cada verbo. Esta labor, realizada siempre con los sujetos más inteligentes y de mayor garantía, se ha llevado a cabo siempre al final de la encuesta, cuando la cantidad de formas verbales recogidas por medios indirectos nos había proporcionado un esqueleto sobre el cual construir nuestro cuestionario morfológico. De este modo hemos procurado comprobar con los mejores sujetos las múltiples formas recogidas de la conversación espontánea de las gentes. Para la sintaxis, difícil de recoger por medio de cuestionario, los ejemplos están tomados casi exclusivamente de la conversación espontánea o dirigida. 

Nos ha permitido emplear estos métodos la calidad excepcional de los sujetos más consultados (cf. § 16). Después de vencer el recelo comprensible de los primeros días, se han sometido durante largas horas a nuestros interrogatorios con una inteligencia y una amabilidad enormes.

Por medio del método descrito hemos tratado de conseguir una estricta comparación de los resultados correspondientes a cada municipio estudiado con los de los otros tres municipios (cf. para el cuestionario, § 15), ya que no teníamos derecho a suponer unidad lingüística en una zona de tan poca unidad geográfica (cf. § 2). 

Ha sido de mucha utilidad en la realización de la encuesta el empleo de un magnetófono portátil. Gracias a él, hemos podido estudiar detenidamente las articulaciones de nuestros sujetos, además de recoger directamente numerosas canciones, romances, relatos varios, etc. (cf. §§ 167-169). Asimismo, se han hecho numerosos dibujos, intercalados en la Quinta Parte (Cosas y Palabras) de este trabajo y una colección de fotografías, reunidas al final. 

\section{\emph{El cuestionario.}}

\subsection{} Ya que no existía ningún trabajo sistemático sobre las hablas de la Montaña, tuvimos que echar mano de los cuestionarios no especializados. Se construyó nuestro cuestionario sobre la base del que empleó M. Alvar para el ALEAr\footnote{\textit{Atlas Lingüístico y Etnográfico de Aragón. Cuestionario,} Sevilla, 1963.}, basado a su vez en el del ALPI. Este cuestionario se ha adaptado y extendido enormemente, quitando aquellos capítulos que no tenían aplicación en nuestra zona (el arado, el carro, muchos cultivos, etc.), ampliando los que quedaban y añadiendo muchos más a base de los vocabularios montañeses publicados (cf. § 12 y Bibliografía). En esta labor de ampliación nos han guiado constantemente ciertas obras de dialectología francesa, sobre todo el cuestionario de P. Nauton para el \textit{Atlas Linguistique et Ethnographique du Massif Central} y el trabajo de Mme. Aub-Büsher sobre el habla de Ranrupt (Bas-Rhin)\footnote{Gertrud Aub-Büsher, \textit{Le parler Rural de Ranrupt. Essai de dialectologie vosgienne}, París, 1962.}.

Sin embargo, el cuestionario que así ha resultado de unos 3500 términos, no ha sido enteramente adecuado a la tarea de investigar todos los aspectos de la cultura de los Montes de Pas. Ha sido necesario un reamoldamiento de muchos capítulos y la adición de varios más, según la experiencia que nos proporcionaba la investigación, adaptación que ha continuado durante casi todo el período de la encuesta. El cuestionario pues, ha sido de carácter muy flexible, aunque de rigidez suficiente como para proporcionar la necesaria base de comparación entre las distintas localidades estudiadas y para evitar que se omitiera de investigar algún aspecto de la vida de cualquier localidad\footnote{Para el cuestionario de morfología, cf. arriba, § 14.}.


\section{\emph{Los sujetos.}}

\subsection{} En cada municipio pedimos a alguna persona conocedora de la zona, fuera alcalde, sacerdote, médico, maestro de escuela, etc., que nos indicara las personas más adecuadas a nuestros fines. Así, durante las primeras semanas de la encuesta, nos dedicamos a interrogar a muchas personas, naturales de los distintos barrios. De este modo pudimos escoger a los dos mejores sujetos de cada sitio y, desde entonces, limitar a ellos nuestra encuesta. En casi todos los casos nos restringimos a personas que habitaban en las zonas más remotas y preferimos siempre a los que tenían la menor cultura y que no habían viajado. Todos nuestros sujetos principales eran mayores de sesenta años y debemos subrayar aquí la alta inteligencia natural de los escogidos. Además de los sujetos principales, consignados abajo, consultamos a varias personas conocedoras de los diversos oficios rurales.

Los sujetos más consultados fueron: 
\begin{enumerate}[a)]
    \item \begin{tabular}{ p{3.5cm} l }
    Vega de Pas. &
    $\left\{
    \begin{tabular}{p{7cm}}
    
    \begin{enumerate}[1)]
       \item Tomás Sañudo Pelayo, campesino, de 70 años; apenas sabe leer, no ha viajado, no ha hecho el servicio militar; natural del barrio de Pandillo, como sus padres (se citará en lo siguiente V3).
       \item Sofía Fernández, esposa del anterior, de 68 años, aproximadamente; no sabe ni leer ni escribir; no ha salido nunca de los Montes de Pas; nacida en Pandillo de padres de la misma localidad (citada: V2). 
       \item Pedro Diego Oria, natural de Yera, de padres nacidos en el municipio; tiene 72 años y era ganadero; no sabe ni leer ni escribir; hizo el servicio militar y ha realizado viajes cortos (citado: V1). 
    \end{enumerate}
    \end{tabular}
    \right.$ \\
    \end{tabular}
    
    \item \begin{tabular}{ p{3.5cm} l }
    San Pedro del Romeral. &
    $\left\{
    \begin{tabular}{p{7cm}}
    
    \begin{enumerate}[1)]
       \item Emilia Ruiz López, nacida en San Pedro como sus padres; tiene 65 años y es viuda; se dedica a sus trabajos domésticos y no sabe leer ni escribir; ha hecho dos visitas a Francia, donde tiene hijos, aparte de lo cual no ha viajado (citada: P1). 
       \item Adolfo Sáinz Revuelta, de 92 años, natural de Bustaleguín, como sus padres; analfabeto; apenas ha viajado (citado: P2).   
    \end{enumerate}
    \end{tabular}
    \right.$ \\
    \end{tabular}

    \item 
    \begin{savenotes}
    \begin{tabular}{ p{3.5cm} l }
    Selaya. &
    $\left\{
    \begin{savenotes}
    \begin{tabular}{p{7cm}}
    
    \begin{enumerate}[1)]
       \item Pilar Gutiérrez, de 75 años, nacida en Cubilla, de padres de Bustantegua; sabe leer y escribir; sirvió como nodriza con dos familias, con las cuales hizo algunos viajes. Mujer de excepcional inteligencia y de rápida comprensión que conoce muy bien el habla local (citada: S1). 
       \item Teresa Fernández, de 80 años, natural de Cubilla, donde nacieron sus padres; no sabe leer ni escribir; nunca ha salido del municipio (citada: S3).
       \item Manuel Ortiz Gutiérrez, de 65 años, nacido en Pisueña; ganadero, que no sabe ni leer ni escribir (citado: S2). 
    \end{enumerate}
    \end{tabular}
    \end{savenotes}
    \right.$
    \end{tabular}
    \end{savenotes}


    \item 
    \begin{savenotes}
    \begin{tabular}{ p{3.5cm} l }
    municipio de San Roque de Riomiera. &
    $\left\{
    \begin{savenotes}
    \begin{tabular}{p{7cm}}
    
    \begin{enumerate}[1)]
       \item José Pérez Ruiz, vecino de Calseca\footnote{Cf. § 13, final.}, donde nació; tiene 55 años y es ganadero; sabe leer un poco (citado: R1). 
       \item Primitivo Crespo, de 68 años, ganadero, nacido en Veolamadera, de padres del municipio; analfabeto, hizo un servicio militar de pocos meses; sin viajar (citado: R2). 
    \end{enumerate}
    \end{tabular}
    \end{savenotes}
    \right.$ \\
    \end{tabular}
    \end{savenotes}
\end{enumerate}

